% =============================================================================
% s8_discussion.tex --- §8 Discussion and Limitations (NEW draft, W7 reorder).
%
% SERVES: the whole paper's honest positioning under the C0-C3 framework ---
%   an analytical / capability paper that charts a reliability x recoverability
%   decision surface and reports honest deployment boundaries, NOT a SOTA chase
%   (STORY 会场 MICCAI 2027 / TMLR, 容诚实负结果).
%
% STORY §8 STRUCTURE (锁定章节顺序 §8.1-§8.4):
%   §8.1 GradProm differentiation (MI lower bound + routing vs gradient repair)
%   §8.2 Why diffusion / generative enhancement fails here (safety, R8)
%   §8.3 the deployable-boundary honest reading (triage global non-superiority)
%   §8.4 Limitations (理论 sketch / 退化 semi-synthetic / q-bar 信号成本 /
%        无真实部署 / E2 contrast 弱 / melanoma 救援缺口未解 / 单中心数据)
%
% SOURCE / MIGRATION:
%   - GradProm段: STORY §"GradProm 对比" weld-shut differentiation (R12).
%   - diffusion段: STORY §8.2 + R8 + main.tex L613-617 (Why generative fails).
%   - boundary段: drafts/s77_dca_triage.tex (sec:exp-triage) global negative +
%     drafts/s71_c0_results.tex / s3_decision_surface.tex C0 three-way verdict.
%   - limitations: STORY §8.4 必含项 + ACCEPTANCE C0/E2/C3.
%
% LOCKED NUMBERS (STORY_FRAMEWORK.md locked table; DO NOT invent / DO NOT alter):
%   C0 decision surface (c0_decision_surface.csv, n=360/cell, verifier PASS):
%     contrast S5 (alpha=0.19) recoverability_delta -0.0355 CI[-0.0783,-0.0013]
%       = the ONLY of 25 cells with significant HURT (enhancement backfires);
%     blur S5 +0.0634 / color_shift S5 +0.0361 = enhancement still HELPs (救得回);
%     completeness all 5 bands cross zero (S5 +0.0236 [-0.0079,+0.0572])
%       = "partially recoverable" trend, statistically NOT significant.
%   E2 per-dim: contrast PSNR 29.11 < degraded 32.29 (enhancement backfires on
%     contrast), color_shift 33.77 < 35 target; brightness/blur PASS.
%   E5 melanoma salvage 5.2% (4/77), net -81 malignant; v6 mask-L1 null 5.2%->5.2%.
%   E6 severe dAUC -0.0559 CI[-0.085,-0.028]. Triage @20%: Direct 0.818 > best
%     variant 0.788 (global honest negative, sec:exp-triage).
%
% RED LINES OBSERVED (STORY R1/R2/R3/R8/R10/R11/R12/R14, drift def. #1,#3,#10):
%   - positioning is CONSTRUCTIVE, never "passive enhancement is unsafe" headline
%     (R14): the contribution is HOW to enhance safely + when to query-for-retake.
%   - NO claim triage globally beats Direct (R3/R11); honest negative stated.
%   - GradProm welded as治标-修梯度 vs our MI-lower-bound + routing (R12), NOT
%     "也修梯度的同类工作".
%   - diffusion ONLY as a rejected cautionary contrast (R8), never our method.
%   - derivation language ("we derive / under assumptions") for the theory, with
%     the R2 forbidden verb avoided; no
%     "universal / always" absolutes (R1); completeness = "partially recoverable"
%     NOT "irreversible".
%   - single-centre / synthetic-degradation framing kept honest; NO BMVC
%     cross-domain rho / ImageNet-C / 5-backbone numbers imported (R10) --- those
%     are BMVC-occupied and only cited (post-notification) elsewhere.
% =============================================================================

\section{Discussion and Limitations}\label{sec:discussion}

% --- §8.1 GradProm differentiation (R12 weld) ----------------------------------
\paragraph{Relation to gradient-repair approaches (GradProm).}
The closest prior observation that enhancement can harm diagnosis is GradProm
\citep{gradprom2025}, which operates in the same ISIC skin-lesion domain. GradProm
treats the symptom by repairing gradients during training (gradient projection /
gradient promotion). Our contribution is orthogonal in both mechanism and scope. We
provide, first, a \emph{diagnosis-preserving information lower bound}: under the stated
assumptions of Lemma~\ref{lem:dp} the feature-level diagnosis-preserving objective
guarantees that the mutual information $I(Z_{\mathrm{enh}};Y)$ retained after
enhancement stays within $\beta\sqrt{\epsilon}$ of the reference (we derive the bound
rather than claim a watertight proof). Second, we provide a \emph{routing decision}
that includes a query-for-retake channel, so that the system can decline to enhance
when enhancement is unsafe rather than always attempting a correction. In short, where
GradProm repairs gradients to mitigate the symptom, we bound the information loss that
enhancement may incur \emph{and} decide when enhancement should not be applied at all.
The two views are complementary rather than competing.

% --- §8.2 Why generative enhancement fails here (R8) ---------------------------
\paragraph{Why generative enhancement is rejected here.}
We deliberately avoid generative restoration (diffusion- or GAN-based) for dermoscopic
enhancement. Generative models optimise perceptual realism, which in dermoscopy can
amount to synthesising plausible-but-absent structure --- pigment networks, vessels, or
globules that the original capture does not contain. Such hallucinated structure is
diagnostically misleading and constitutes a medical-safety hazard, since the
downstream reader or classifier cannot distinguish invented detail from real
morphology. We therefore use deterministic restoration with a quality-conditional
modulation and a diagnosis-preserving objective, trading some perceptual sharpness for
the information-preservation property analysed in \cref{sec:prop3lemma3}. This
trade-off is the design intent, not a shortcoming: a system that is permitted to
invent lesion structure cannot offer the diagnosis-preservation guarantee that
motivates the present work. Diffusion-based restoration thus enters our discussion
only as a rejected cautionary contrast, never as a component of the method.

% --- §8.3 the deployable-boundary honest reading -------------------------------
\paragraph{The honest reading of the deployable boundary.}
Our decision-curve and triage analysis (\cref{sec:exp-triage}) supports a deliberately
narrow positive claim and an explicit negative one. The negative claim is that, at a
fixed referral threshold, globally applied quality-aware triage \emph{does not}
outperform direct diagnosis: on the worst-case low-quality subset direct diagnosis
attains the highest auto-handled sensitivity ($0.818$), the strongest tested bottleneck
variant only matches it ($0.788$), and the maximum-net-benefit confidence intervals of
all methods overlap ($0.179$--$0.192$). We make no claim that triage raises net benefit
globally, and we do not assert that the agent's global risk is at most that of direct
diagnosis. The positive claim is correspondingly local: \emph{within} the
moderate-degradation window, the diagnosis-preserving enhancer does not lose
discriminative information (\cref{sec:exp-preserve,sec:exp-dploss}), and the
reliability $\times$ recoverability decision surface (\cref{sec:exp-surface}) shows that
recoverability is dimension- and severity-dependent rather than uniform. Reading that
surface honestly yields a three-way verdict: some degradations remain recoverable even
at the most severe level tested (blur and colour-shift enhancement still improve
diagnostic AUC at the strongest setting, with recoverability deltas whose CIs exclude
zero); one degradation backfires (the most extreme contrast reduction, $\alpha=0.19$,
is the single cell of twenty-five where enhancement significantly lowers AUC,
recoverability delta $-0.0355$, $95\%$ CI $[-0.0783,-0.0013]$); and one degradation is
only \emph{partially recoverable} with a trend that does not reach significance
(field-of-view truncation / completeness, where every severity band's recoverability
CI crosses zero). This partitioning --- recover, refuse-and-retake, or proceed with
caution --- is precisely the decision the query-for-retake gate is built to make, and
it is why the contribution is framed as charting and bounding that decision rather than
claiming a universal enhancement win.

% --- §8.4 Limitations ----------------------------------------------------------
\paragraph{Limitations.}
We state the limitations of this work plainly, as several of them are also the
motivation for the query-for-retake gate.
(1)~\textbf{Theory as sketches under assumptions.} The theoretical results
(Lemma~\ref{lem:dp}, Proposition~\ref{prop:entropy}, Theorem~\ref{thm:agent-compact}) are
derived under explicitly stated assumptions and are arguments for why the bounds hold
rather than fully watertight proofs in every regime; in particular,
Theorem~\ref{thm:agent-compact} is a band-local conditional bound, and its
within-band guarantee does not aggregate into a global ordering
(Remark~\ref{rem:no-global}).
(2)~\textbf{Semi-synthetic degradation.} The benchmark applies controlled synthetic
degradations along five physical axes; real consumer-capture statistics may differ
from this protocol, and prospective validation on genuinely degraded clinical captures
remains open.
(3)~\textbf{Quality-signal cost.} The routing policy depends on a per-input quality
scalar; learning a reliable quality channel adds annotation and training cost, and
re-deployment within dermatology to a markedly different capture distribution would
require re-fitting that channel (we restrict our quality-signal claims to the
dermatology modality and do not assert cross-modality universality).
(4)~\textbf{No real-world deployment validation.} We report no prospective clinical
deployment and no clinician confirmation; all clinical statements rest on
decision-curve analysis and triage simulation
(\cref{sec:exp-triage}, Appendix~\ref{app:a20}), with their stated limits. The system
is a research framework, not a medical device.
(5)~\textbf{Single-centre, public data only.} Evaluation uses publicly available
de-identified dermatology datasets; the diagnosis-preservation conclusions are
additionally measured against a single frozen diagnostic probe, so they are
probe-dependent, and a second independent diagnostic backbone and multi-centre data
would strengthen them.
(6)~\textbf{Per-dimension enhancement is uneven.} Enhancement does not improve every
degradation axis: on the contrast axis it backfires (enhanced PSNR $29.11$~dB falls
below the degraded input's $32.29$~dB), and colour-shift restoration does not reach our
quality target ($33.77$~dB $<35$~dB), so these axes are honest failure modes that
trigger the retake / caution channels rather than enhancement.
(7)~\textbf{The malignant salvage gap is unresolved.} Enhancement salvages only
$5.2\%$ ($4/77$) of misclassified melanomas while damaging $31.0\%$ of correctly
classified ones (a net loss of $81$ malignant cases), and a mask-weighted
reconstruction retrain leaves this unchanged ($5.2\%\to5.2\%$, an honest null). Closing
this malignant-rescue gap --- for instance through diagnosis-aware or adversarial
reconstruction objectives, or through learned rather than fixed routing thresholds ---
is the central item of future work. We report these negative results rather than omit
them, because they are the genuine evidence that a query-for-retake gate is a necessary
safety mechanism and not a project shortcoming.
(8)~\textbf{Uneven cross-source transfer and a skin-tone disparity.} A supplementary
analysis (Appendix~\ref{app:gen-fair}) shows that transfer across dermatology acquisition
sources is highly uneven --- strong near the training distribution (HAM10000) but
degrading towards chance on the most distant source (DermNet) --- and that diagnostic AUC
is systematically lowest on darker (Fitzpatrick V--VI) skin across every method examined.
We surface both as honest boundaries on the system's reach rather than resolved properties.

\paragraph{Future work.}
Two directions follow directly from the honest boundary. First, replacing the
fixed routing thresholds with \emph{learned} thresholds is the most promising route to
turning the band-local guarantee (\cref{sec:exp-triage}, Remark~\ref{rem:no-global})
into a global utility improvement, since the global triage negative result is obtained
under fixed thresholds. Second, closing the malignant salvage gap calls for
reconstruction objectives that reach the diagnostic decision boundary rather than
per-pixel fidelity alone. A formal cost-effectiveness study with jurisdiction-specific
inputs (Appendix~\ref{app:a20}) and prospective clinical validation are necessary
before any deployment claim.
